\section[光滑性]{光滑性（12 月 18 日）}

\srcpage{79}

\subsection{上次的回顾}

我们简要回顾了菱形的 étale 上同调。设 $\Lambda$ 是被某个与 $p$ 互素的整数 $n$ 湮灭的环。对任意小 v-叠，我们定义了 $D_{\et}(X,\Lambda)$。一般情形下它由下降定义，但若 $X$ 是满足 $\dim.\mathrm{tr.g.}(X)<\infty$ 的局部空间菱形，则它与 $D(X_{\et},\Lambda)$ 一致。它是一个（闭的）对称幺半三角范畴。

对任意映射 $f\colon Y\to X$，我们有 $f^*$ 与 $Rf_*$。

若 $f$ 在局部空间菱形中可表、可紧化，且局部地 $\dim.\mathrm{tr.g.}(f)<\infty$，则有函子
\[
Rf_!\colon D_{\et}(Y,\Lambda)\to D_{\et}(X,\Lambda),
\]
带右伴随 $Rf^!$。

这满足「六函子形式化」的一切性质。实践中困难的部分是理解 $Rf^!$。

\subsection{为什么要求局部空间菱形？}

上次我们没有解释的一件事——它也是「局部空间」这一定义的主要动机之一——是我们为什么要求映射在局部空间菱形中可表。要点在于我们希望 $Rf_!$ 与一切直和交换。这与「有右伴随」等价，而由伴随函子定理，后者又本质上等价于与任意直和交换。最终，这些假设将保证这一点。

\begin{example}
我们总在有界导出范畴上考虑这些函子。若 $A_n\in D_{\et}(Y,\Lambda)$（$n\ge 0$）集中在次数 $0$，则 $\bigoplus A_n[n]\cong\prod A_n[n]$（由左完备性——自然映射是否为同构，可通过察看上同调层来检验）。于是
\[
Rf_!\Bigl(\bigoplus A_n[n]\Bigr)\cong Rf_!\Bigl(\prod_n A_n[n]\Bigr)\cong\prod_n\bigl(Rf_!A_n\bigr)[n].
\]
我们希望它与 $\bigoplus_n(Rf_!A_n)[n]$ 同构。这只有在 $Rf_!$ 具有有限上同调维数时才可能成立。
\end{example}

于是有一般原则：一个函子（在三角左完备范畴上）只有当其具有有限上同调维数时，才可能与一切直和交换。

一个关键的要素是：

\begin{theorem}[Scheiderer '94]
若 $T$ 是任意谱拓扑空间，则 $T$ 的上同调维数不大于 $T$ 的 Krull 维数。
\end{theorem}

\begin{remark}
在诺特性假设下，该定理早得多就由 Grothendieck 证明了。这是一个引人注目的定理，因为 Krull 维数是纯局部的不变量，而上同调维数是整体的量。该结论对紧 Hausdorff 空间严重失效（后者的 Krull 维数总是 $0$）。
\end{remark}

我们所施加的有限性通过「几何超越维数」$\dim.\mathrm{tr.g.}$ 实现，它是纯粹逐点的条件。它足以侦测 Krull 维数；但要控制上同调维数，我们需要空间是（局部）谱的，而这正是「局部空间」所提供的东西。总之，我们需要「局部空间」假设来控制上同调维数。

\srcpage{80}

\subsection{上同调光滑性}

\begin{definition}
固定 $\ell\ne p$。则上述映射 $f\colon Y\to X$（可紧化、在局部空间菱形中可表、且局部地 $\dim.\mathrm{tr.g.}(f)<\infty$）称为 \emph{$\ell$-上同调光滑}的，若经任意基变换后，有函子的同构
\[
Rf^!\cong\mathbb{D}_f\otimes f^*\colon D_{\et}(X,\mathbb{F}_\ell)\to D_{\et}(Y,\mathbb{F}_\ell),
\]
其中 $\mathbb{D}_f$ 局部同构于 $\mathbb{F}_\ell[n]$（某个 $n\in\ZZ$）。
\end{definition}

\begin{remark}
若此条件满足，则 $\mathbb{D}_f=Rf^!\mathbb{F}_\ell$，且它的构造与任意基变换交换。

反过来（这多少有些出人意料），若 $Rf^!\mathbb{F}_\ell$ 可逆，且它的构造沿 $Y\to X$ 与基变换交换，则 $f$ 是 $\ell$-上同调光滑的。
\end{remark}

\begin{definition}
称 $f\colon Y\to X$ 是上同调光滑的，若它对一切 $\ell\ne p$ 都是 $\ell$-上同调光滑的。此时对任意 $\Lambda$，$Rf^!\Lambda$ 局部同构于 $\Lambda[n]$。
\end{definition}

\begin{example}
回忆 $\mathbb{B}_C$ 是闭单位圆盘。映射 $\mathbb{B}_C^\diamond\to(\Spa C)^\diamond$ 是上同调光滑的。这可由 Huber 的结果推出，他研究过 $\mathbb{B}_C$ 这类空间的上同调。先验地看，我们需要对任意基变换作检验，这可能很困难。最初的论证是通过基变换到严格完全间断空间。不过，你也可以使用注 15.5 中的反方向。
\end{example}

\begin{example}
若 $f\colon Y\to X$ 是 $\Zp$ 上解析 adic 空间之间的光滑映射（即局部 étale 于某个球之上），则 $f^\diamond\colon Y^\diamond\to X^\diamond$ 是上同调光滑的。
\end{example}

\begin{example}
设 $E$ 是 $\Qp$ 的有限扩张。则 $(\Spa E)^\diamond\to(\Spa\mathbb{F}_q)^\diamond$ 是上同调光滑的。更一般地，回忆 $\Div^d=(\Div^1)^d/S_d$。则 $\Div^d\to*$ 是上同调光滑的。

事实上，这对非解析空间也成立：$(\Spa\cO_E)^\diamond\to(\Spa\mathbb{F}_q)^\diamond$ 是上同调光滑的。回忆 $\Spa\cO_E$ 与 $\Spa\mathbb{F}_q$ 都不是解析 adic 空间，故 $(\Spa\cO_E)^\diamond$ 与 $(\Spa\mathbb{F}_q)^\diamond$ 不是菱形。但它们是小的 v-层，因此我们的形式化仍然适用。
\end{example}

\begin{example}
若 $f\colon Y\to X$ 上同调光滑，且 pro-$p$ 群 $G$ 自由地作用在 $Y/X$ 上，则 $f/G\colon Y/G\to X$ 仍然上同调光滑。其逆命题则大错特错！若 $G$ 无限，则映射
\[
Y=\Spa C\times\ul{G}\to\Spa C=X
\]
不是上同调光滑的。
\end{example}

\begin{example}
考虑开的 Schubert 胞腔 $\Gr_{G,\mu}$。映射 $\Gr_{G,\mu}\to(\Spa E)^\diamond$ 上同调光滑，但 $\Gr_{G,\le\mu}$ 一般不是。
\end{example}

\begin{example}
若 $\mathcal{E}$ 是 $X_S$ 上的向量丛，且其所有纤维只有正的 HN 斜率，则 $BC(\mathcal{E})\to S$ 上同调光滑。若斜率 $\ge 0$ 则不然：例如 $\mathcal{E}=\cO_{X_S}$ 时 $BC(\mathcal{E})=\ul{E}$，而它并非上同调光滑。
\end{example}

\begin{example}
映射 $f\colon Y\to X$ 上同调光滑这一条件，可在 $X$ 上 v-局部地检验，前提是 $f$ 可紧化、在局部空间菱形中可表、且局部地 $\dim.\mathrm{tr.g.}(f)<\infty$。
\end{example}

所有这些例子都是从 $\mathbb{B}_C$ 出发建立的，而这本质上已包含在 Huber 的工作中。随后你观察到：上同调光滑性在取 pro-$p$ 群之商下是稳定的，因为这类群的 $\ell$-上同调维数为 $0$。

\srcpage{81}

\subsection{Artin 叠}

类比：
\[
\begin{tikzcd}
\text{概形} \arrow[r, leftrightarrow] & \text{代数空间} \arrow[r, leftrightarrow] & \text{Artin 叠} \\
\text{完满空间} \arrow[r, leftrightarrow] & \text{局部空间菱形} \arrow[r, leftrightarrow] & \text{Artin v-叠}
\end{tikzcd}
\]

\begin{definition}
称小 v-叠 $X$ 是 \emph{Artin} 的，若
\begin{itemize}
\item $\Delta_X\colon X\to X\times X$ 在局部空间菱形中可表；
\item 存在上同调光滑的满射 $f\colon Y\to X$，其中 $Y$ 是局部空间菱形。
\end{itemize}
\end{definition}

对 Artin 叠 $X$，可以定义 $X\to*$ 是「上同调光滑的」意味着什么：它指的是，对定义 15.14 中的 $f\colon Y\to X$，复合 $Y\to X\to*$ 上同调光滑。我们可以把对偶复形的概念推广到这样的 $X$，并通过对偶复形的次数得到维数的概念。（先验地，这依赖于 $\ell$。）

\begin{theorem}
$\Bun_G$ 是上同调光滑的 Artin v-叠（维数为 $0$，意即：对 $\pi\colon\Bun_G\to*$，$R\pi^!\Lambda$ 局部同构于 $\Lambda[0]$）。
\end{theorem}

\begin{proof}[证明梗概]
要证 Beauville--Laszlo 单值化
\[
\pi\colon\Gr_{G,\mu}/G(E)\to\Bun_G
\]
是维数为 $\langle2\rho,\mu\rangle$ 的上同调光滑映射，且 $\Gr_{G,\mu}/G(E)$ 也是维数为 $\langle2\rho,\mu\rangle$ 的上同调光滑空间。

$\pi$ 的纤维是 $\Gr_{G,\mu^{-1}}$ 中的开集。（纤维是某个丛的那些为平凡的修正，而我们已经看到「平凡」是开条件。）

我们断言：qcqs 的上同调光滑映射是开映射，因而该映射的（开）像具有所要的性质。为论证此断言，设 $f\colon Y\to X$ 上同调光滑且 qcqs。则 $Rf_!$ 保持可构造层：它保持紧对象，因为其右伴随 $Rf^!$ 保持直和，而紧对象正是可构造层。于是 $Rf_!\mathbb{F}_\ell$ 可构造，从而其支撑是开的且拟紧。

取这样的图表之并 $\bigcup_\mu$，由定理 12.25，这覆盖了 $\Bun_G$。
\end{proof}

\begin{example}
特别地，$[\mathrm{pt}/G(E)]\cong\Bun^1_G\subset\Bun_G$ 在 $*$ 上是上同调光滑的。
\end{example}

\subsection{$\Bun_G$ 的更好图表}

为研究 $\Bun_G$，我们需要更好的光滑图表。

\begin{example}
设 $G=\mathrm{GL}_2$。我们有特殊化 $\cO(1/2)\rightsquigarrow\cO\oplus\cO(1)$。这两个点构成开子叠 $U\subset\Bun_G$。

我们想要 $U$ 的一个好的图册。设 $b\in B(G)$ 是对应于 $\cO\oplus\cO(1)$ 的等晶体。设 $\mathcal M_b$ 是扩张
\[
0\to\mathcal L\to\mathcal E\to\mathcal L'\to0
\]
的模空间，其中 $\mathcal L$ 是次数 $0$ 的线丛，$\mathcal L'$ 是次数 $1$ 的线丛。

于是得到映射
\[
\pi_b\colon\mathcal M_b\to\Bun_{\mathrm{GL}_2}=\Bun_G,
\]
把 $(\mathcal L\to\mathcal E\to\mathcal L')$ 送到 $\mathcal E$。
\end{example}

\begin{theorem}
映射 $\pi_b$ 上同调光滑。
\end{theorem}

\srcpage{82}

$\mathcal M_b$ 的结构如何？我们有 $\mathcal M_b=\wt{\mathcal M}_b/(E^\times\times E^\times)$，其中 $\wt{\mathcal M}_b$ 参数化
\[
0\to\cO\to\mathcal E\to\cO(1)\to0.
\]
于是 $\wt{\mathcal M}_b=BC(\cO(-1)[1])$ 是负 Banach--Colmez 空间。

我们有一个图表
\[
\begin{tikzcd}
\wt{\mathcal M}_b^\circ & \wt{\mathcal M}_b \arrow[l] & * \arrow[l] \arrow[d] \\
\Bun^{b'}_G \arrow[u] & U \arrow[l] \arrow[u] & \Bun^b_G \arrow[l] \arrow[u]
\end{tikzcd}
\]
其中 $b'\leftrightarrow\cO(1/2)$。

这是因为 $\cO(1)$ 的每个非分裂扩张都同构于 $\cO(1/2)$。我们有 $\wt{\mathcal M}_b^\circ=BC(\cO(-1)[1])\setminus\{0\}$。它可以显式地表示为
\[
(\Spa k((t)))^\diamond/\mathrm{SL}_1(D_{1/2}),
\]
其中 $D_{1/2}$ 是 $E$ 上的四元数代数。事实上，$\mathrm{SL}_1(D_{1/2})$ 由 $\cO(1/2)$ 的这样一些自同构组成：它们保持其相伴分次的给定平凡化。在选定与这个标准扩张的同构之后，剩下的只是给出 $\cO(1/2)$ 的一个非零截面，而后者由 Banach--Colmez 空间 $BC(\cO(1/2))\setminus0$ 参数化，即一个去心开单位圆盘。

一般地，我们将定义上同调光滑映射
\[
\pi_b\colon\mathcal M_b\to\Bun_G.
\]

注意 $[\mathrm{pt}/G_b(E)]\xrightarrow[\text{上同调光滑}]{}\Bun^b_G\to\Bun_G$。所以你也许会问：能否把上同调光滑映射 $[\mathrm{pt}/G_b(E)]\to\Bun^b_G$ 延拓到一个小邻域？$\mathcal M_b$ 所做的事情正是这个。你可以把它看作「$\Bun_G$ 在 $[\mathrm{pt}/G_b(E)]\to\Bun^b_G$ 处的严格 Hensel 化」。它嵌入图表
\[
\begin{tikzcd}
\wt{\mathcal M}_b \arrow[r] & \mathcal M_b \\
* \arrow[u, "s"] \arrow[r] & {[\mathrm{pt}/G_b(E)]} \arrow[u, "s"]
\end{tikzcd}
\]
其中截面 $s$ 满射到 $\Bun^b_G\subset\Bun_G$ 的原像上，且 $\wt{\mathcal M}^\circ_b:=\wt{\mathcal M}_b\setminus*\to\Bun^{>b}_G$。

$\wt{\mathcal M}^\circ_b$ 是空间菱形。

\begin{warning}
注意映射 $\wt{\mathcal M}_b\to*$ 在局部空间菱形中可表，但 $\wt{\mathcal M}_b$ 本身不是局部空间菱形，因为它不是解析的。
\end{warning}

\begin{example}
考虑 $\Spa k[[t]]$，它表示 $\Spa(R,R^+)\mapsto R^{\circ\circ}$。映射 $\Spa k[[t]]\to*$ 在局部空间菱形中可表，因为它到任意 $S$ 的基变换是 $S$ 上的开单位圆盘。但 $\Spa k[[t]]$ 不是菱形，因为它有非解析点。

然而 $\Spa k[[t]]\setminus\Spa k=\Spa k((t))$ 是空间菱形。
\end{example}

因此我们需要注意：一个抽象性质与它在 $*$ 上的相对性质可以非常不同。$\Spa k((t))$ 到 $*$ 的映射不拟紧，而 $\Spa k[[t]]$ 拟紧。

事实上，这一现象对论证至关重要。归根结底，对我们关键的是证明有限性（例如 shtuka 模空间的上同调群的有限性）。能够改变视角、使某物变得拟紧，这套论证的技术性工作很大一部分就在这里。

\srcpage{83}

对 $\mathrm{GL}_n$，等晶体 $b\in B(\mathrm{GL}_n)$ 对应于一个 $\mathbb{Q}$-分次向量丛 $\mathcal E=\bigoplus_{\lambda\in\mathbb{Q}}\mathcal E^\lambda$，其中 $\mathcal E^\lambda$ 是斜率为 $\lambda$ 的半稳定对象。于是 $\mathcal M_b$ 参数化诸 $\mathcal E^\lambda$ 的迭代扩张。

\begin{definition}
设 $\mathcal M$ 是把 $S\in\Perf_{\Fq}$ 送到正合 $\otimes$-函子
\[
\Rep_E(G)\to\underbrace{\{\mathbb{Q}\text{-滤过向量丛 }\mathcal E\supset\mathcal E^{\le\lambda}\}}_{\text{正合范畴}}
\]
的模空间，使得 $\mathcal E^\lambda=\mathcal E^{\le\lambda}\big/\bigcup_{\lambda'<\lambda}\mathcal E^{\le\lambda'}$ 是斜率为 $\lambda$ 的半稳定对象。
\end{definition}

\begin{warning}
注意这个滤过与 HN 滤过「相反」！（斜率递增。）
\end{warning}

我们有一个映射 $\mathcal M\to\Bun_G$，把 $(\mathcal E\supset\mathcal E^{\le\lambda})$ 送到 $\mathcal E$。

$\mathcal M$ 上还有一个映射 $\mathcal M\to\coprod_{b\in B(G)}[\mathrm{pt}/G_b(E)]$，把 $(\mathcal E\supset\mathcal E^{\le\lambda})$ 送到 $\bigoplus\mathcal E^\lambda$。这是一个取值于 $\mathbb{Q}$-分次向量丛的正合 $\otimes$-函子，其中 $\mathcal E^\lambda$ 是斜率为 $\lambda$ 的半稳定对象。这个目标范畴等价于等晶体的范畴，后者又等价于 $[\mathrm{pt}/G_b(E)]$。这诱导出 $\mathcal M=\bigcup_{b\in B(G)}\mathcal M_b$，其中 $\mathcal M_b$ 是 $b$ 上的纤维。

\begin{theorem}
映射 $\pi_b\colon\mathcal M_b\to\Bun_G$ 上同调光滑。
\end{theorem}

\begin{remark}
对 $G=\mathrm{GL}_n$，这可直接用正面攻击证明。对一般的 $G$ 则相当困难，我们将需要引入光滑性的「Jacobian 判别法」。
\end{remark}
